% ---------------------------------------------------------------------
%  08-metrics.tex -- Metrics that fall out of the graph
%  The formulas (labels eq:leadtime ... eq:unprop) live in
%  appendix-a-formal.tex, subsection app:metrics. % F10 -- Metrics derived from the graph, computed from the Condition A example
% (figure "software-line"). Illustrative data. Requires hpb.sty v1.0 and
% pgfplots (compat=1.18). Ordinal columns are those of fig:software-line.
\begin{figure}[t]
\centering
\resizebox{\linewidth}{!}{%
\begin{tikzpicture}[
    font=\sffamily\fontsize{7}{8.4}\selectfont,
    mx/bar/.style  = {draw=hpbInk, line width=0.8pt, rounded corners=2pt},
    mx/in/.style   = {font=\sffamily\fontsize{6.5}{7.8}\selectfont, text=hpbInk,
                      inner sep=1pt},
    mx/ann/.style  = {hpb/callout, align=left},
  ]
  \pgfplotsset{
    mx/common/.style = {
      width=13.2cm, scale only axis, xmin=0, xmax=17,
      axis line style={draw=hpbDiscInk, line width=0.6pt},
      tick style={draw=hpbDiscInk, line width=0.6pt},
      tick label style={font=\sffamily\fontsize{7}{8.4}\selectfont, text=hpbInk},
      label style={font=\sffamily\fontsize{7}{8.4}\selectfont, text=hpbInk},
      xtick={1,5,10,16},
      xmajorgrids, major x grid style={draw=hpbDiscInk, line width=0.5pt,
        dash pattern=on 2pt off 2pt},
      clip=false, axis on top=false,
    },
  }
  % ------------------------------------------------------------------ (a)
  \begin{axis}[mx/common, name=wip, height=3.3cm,
      ymin=0, ymax=4.3, ytick={0,1,2,3}, xticklabels={},
      ymajorgrids, major y grid style={draw=hpbGuide, line width=0.4pt},
      ylabel={open branches $|W_\theta(t)|$},
      stack plots=y, const plot mark left,
      every axis plot/.append style={draw=hpbInk, line width=0.8pt, mark=none},
      area legend,
      reverse legend,
      legend style={draw=none, fill=white, fill opacity=0.85, text opacity=1,
        at={(0.008,0.975)}, anchor=north west, inner sep=1pt,
        legend cell align=left, row sep=-1pt,
        font=\sffamily\fontsize{7}{8.4}\selectfont, text=hpbInk},
    ]
    % stacked bottom-to-top as in the diagram: T (lane -2), X (lane -1), I (lane +1)
    \addplot[fill=hpbTask] coordinates {
      (0,0)(1,0)(2,0)(3,0)(4,0)(5,0)(6,0)(7,1)(8,1)(9,1)(10,1)(11,1)(12,1)
      (13,1)(14,1)(15,1)(16,0)(17,0)} \closedcycle;
    \addplot[fill=hpbIss] coordinates {
      (0,0)(1,0)(2,1)(3,1)(4,1)(5,0)(6,1)(7,1)(8,1)(9,1)(10,0)(11,0)(12,1)
      (13,1)(14,1)(15,1)(16,0)(17,0)} \closedcycle;
    \addplot[fill=hpbImp] coordinates {
      (0,0)(1,0)(2,0)(3,1)(4,1)(5,1)(6,1)(7,1)(8,1)(9,1)(10,0)(11,1)(12,1)
      (13,1)(14,1)(15,1)(16,0)(17,0)} \closedcycle;
    \legend{Task (T), Issue (X), Improvement (I)}
    % type letters inside the bands: the stack is readable without color
    \foreach \c/\y/\l in {12.5/0.5/T, 3/0.5/X, 12.5/1.5/X,
                          4/1.5/I, 12.5/2.5/I}{
      \edef\mxtmp{\noexpand\node[font=\noexpand\sffamily\noexpand\bfseries
        \noexpand\fontsize{7}{8.4}\noexpand\selectfont, text=hpbInk,
        inner sep=0pt] at (axis cs:\c,\y) {\l};}\mxtmp
    }
    % baseline tags above the panel
    \foreach \c/\v in {1/v1.0, 5/v1.0.1, 10/v1.1, 16/v1.2}{
      \edef\mxtmp{\noexpand\node[hpb/tag, anchor=south, yshift=3pt]
        at (axis cs:\c,4.3) {\v};}\mxtmp
    }
    % annotations
    \draw[hpb/leader] (axis cs:9.7,2.2) |- (axis cs:8.6,3.72)
      node[mx/ann, anchor=east]
      {release of v1.1: I-21 and X-102\\close together ($|C|=2$)};
    \draw[hpb/leader] (axis cs:15.7,2.2) |- (axis cs:15.1,3.72)
      node[mx/ann, anchor=east]
      {batch merge at v1.2:\\three branches close ($|C|=3$)};
  \end{axis}
  % ------------------------------------------------------------------ (b)
  \begin{axis}[mx/common, at={(wip.south west)}, anchor=north west,
      yshift=-0.55cm, height=3.6cm,
      ymin=0.4, ymax=6.6, ytick={1,...,6}, y dir=reverse,
      yticklabels={I-21, I-22, X-101, X-102, X-103, T-8},
      y tick label style={font=\sffamily\bfseries\fontsize{7}{8.4}\selectfont},
      major y tick style={draw=none},
      xticklabels={Jan, Feb, Apr, Jul},
      xlabel={ordinal time (one column per event; calendar month of each baseline)},
      ylabel={branch},
    ]
    % one bar per branch: from opening t_o(b) (start node) to closing t_c(b)
    % (release of the merge baseline);  #1 row, #2 type color, #3 letter,
    % #4 t_o column, #5 t_c column, #6 fork -> merge baseline, #7 D(b)
    \foreach \r/\tc/\tl/\o/\e/\fm/\dd in {
        1/hpbImp/I/3/10/{v1.0 $\to$ v1.1}/2,
        2/hpbImp/I/11/16/{v1.1 $\to$ v1.2}/1,
        3/hpbIss/X/2/5/{v1.0 $\to$ v1.0.1}/1,
        4/hpbIss/X/6/10/{v1.0.1 $\to$ v1.1}/1,
        5/hpbIss/X/12/16/{v1.1 $\to$ v1.2}/1,
        6/hpbTask/T/7/16/{v1.0.1 $\to$ v1.2}/2}{
      \edef\mxtmp{%
        \noexpand\draw[mx/bar, fill=\tc!55]
          (axis cs:\o,\r-0.27) rectangle (axis cs:\e,\r+0.27);
        \noexpand\node[hpb/event, fill=\tc] at (axis cs:\o,\r) {\tl};
        \noexpand\node[mx/in, anchor=west] at (axis cs:\o,\r) [xshift=9pt]
          {\unexpanded\expandafter{\fm}};
        \noexpand\node[mx/in, font=\noexpand\sffamily\noexpand\bfseries
          \noexpand\fontsize{7}{8.4}\noexpand\selectfont]
          at (axis cs:17.55,\r) {\dd};}%
      \mxtmp
    }
    \node[mx/in, font=\sffamily\bfseries\fontsize{7}{8.4}\selectfont]
      at (axis cs:17.55,0.4) [anchor=south, yshift=1pt] {$D(b)$};
  \end{axis}
  % panel letters
  \node[font=\sffamily\bfseries\fontsize{7.5}{9}\selectfont, anchor=north east]
    at ([xshift=-1.15cm]wip.north west) {(a)};
  \node[font=\sffamily\bfseries\fontsize{7.5}{9}\selectfont, anchor=north east]
    at ([xshift=-1.15cm,yshift=-0.55cm]wip.south west) {(b)};
\end{tikzpicture}}
\caption{Flow metrics read straight from the graph, for the data of
situation~A (\cref{fig:software-line}) on the same ordinal columns. (a)~Work
in progress per type: each opening raises a step, each closure lowers it, and
a release drops it by its batch size $|C(v)|$. (b)~One bar per branch, from
its opening to the release that closed it, with its fork and merge versions
and its merge distance $D(b)$. Lead times in calendar time come from the
event dates, not from the length of the bars.}
\label{fig:metrics}
\end{figure}

%  appears here exactly once.
% ---------------------------------------------------------------------
\section{Metrics That Fall Out of the Graph}
\label{sec:metrics}

A useful side effect of the rules is that measurement comes for free.
Every branch has a type, an origin baseline and an opening date and, once
closed, exactly one terminal with a closing date. That is enough to compute
the flow and quality metrics of \cref{tab:metrics} straight from the
graph, with no extra data collection and no spreadsheet kept on the side.
The exact definitions are in \cref{app:metrics}.\footnote{Two details keep
the counts honest. Sub-tasks are steps, not branches (\cref{rule:typed}),
so every count counts a unit of change once. And a reopened branch is rare
but possible: each period during which a branch is open or targeted is a
\emph{visit}, lead times are measured per visit, and reopens are reported
separately.}

All of them are \emph{descriptive}. They exist to bring good questions to
the meetings of \cref{subsec:method-cadence}, not to rank teams or people
(\cref{subsec:metrics-goodhart}).

\begin{table}[tbp]
  \centering
  \small
  \caption{The HPB metrics: the question each one answers, how to read it,
    and the meeting it feeds (\cref{subsec:method-cadence}). Definitions in
    \cref{app:metrics}.}
  \label{tab:metrics}
  \renewcommand{\arraystretch}{1.1}
  \begin{tabularx}{\linewidth}{@{}>{\raggedright\arraybackslash}p{2.3cm}
      >{\raggedright\arraybackslash}X
      >{\raggedright\arraybackslash}X
      >{\raggedright\arraybackslash}p{1.9cm}@{}}
    \toprule
    \textbf{Metric} & \textbf{Question it answers} & \textbf{How to read it}
      & \textbf{Feeds} \\
    \midrule
    Lead time
      & How long from accepting a request to having it in a published
        version (or deciding not to)?
      & Merged and discarded apart, by type. Ends at publication: the
        time customers feel.
      & Line review \\
    Age of open work
      & What has been open too long?
      & Older than the usual lead time of its type: on the triage agenda.
      & Triage \\
    Work in progress
      & How much is open now, per type?
      & Open issues rising and not brought down by releases: an early sign
        of a quality problem.
      & Triage, planning \\
    Throughput
      & How much do we deliver per period?
      & Use windows that span several releases; read it with lead time.
      & Line review \\
    Merge and discard ratios
      & Of what we closed, how much did we deliver?
      & Only by discard reason. Many duplicates: intake problem. Many
        \emph{not viable} lab branches: often healthy.
      & Line review \\
    Merge distance
      & How many releases did a change wait?
      & Large values: work pushed back again and again; its specification
        may be stale.
      & Planning \\
    Release batch size
      & How much does a version contain?
      & Compare a planned set with recent releases before committing it.
      & Planning \\
    Escaped defects
      & Which issues were seen in each version, per change shipped?
      & Compare only at similar exposure; ``seen in'' is not ``introduced
        by''.
      & Line review \\
    Gate slip
      & How late did the version ship?
      & Never alone: with the branches moved out and discarded from it.
      & Gate review \\
    Maintenance load
      & How much new work still lands on an old reference?
      & Sets life-cycle dates with numbers instead of opinions.
      & Line review \\
    Unported fixes
      & Which fixes affect another maintained stream and were not ported?
      & Should be zero after every line review.
      & Line review \\
    \bottomrule
  \end{tabularx}
\end{table}

\paragraph{Flow.}
Lead time ends when the version is \emph{published}, not when engineering
is done: it measures how long until the change is available in the
product, the same idea as lead time in software delivery
\citep{forsgren2018accelerate}, applied to firmware and hardware too. Its
leading counterpart is the age of what is still open. Work in progress is
a staircase over time, one step up at each fork, one down at each terminal
and a big drop at each release; \cref{fig:metrics} shows it for
situation~A. Little's law \citep{little1961} ties work in progress,
throughput and lead time together. I use it as a sanity check (if the
three numbers do not roughly agree, the window is too short or the system
is not stable) and as a planning lever, in the spirit of
\citet{reinertsen2009flow}: at a given delivery rate, the only way to
shorten lead time is to have less open work.

% F10 -- Metrics derived from the graph, computed from the Condition A example
% (figure "software-line"). Illustrative data. Requires hpb.sty v1.0 and
% pgfplots (compat=1.18). Ordinal columns are those of fig:software-line.
\begin{figure}[t]
\centering
\resizebox{\linewidth}{!}{%
\begin{tikzpicture}[
    font=\sffamily\fontsize{7}{8.4}\selectfont,
    mx/bar/.style  = {draw=hpbInk, line width=0.8pt, rounded corners=2pt},
    mx/in/.style   = {font=\sffamily\fontsize{6.5}{7.8}\selectfont, text=hpbInk,
                      inner sep=1pt},
    mx/ann/.style  = {hpb/callout, align=left},
  ]
  \pgfplotsset{
    mx/common/.style = {
      width=13.2cm, scale only axis, xmin=0, xmax=17,
      axis line style={draw=hpbDiscInk, line width=0.6pt},
      tick style={draw=hpbDiscInk, line width=0.6pt},
      tick label style={font=\sffamily\fontsize{7}{8.4}\selectfont, text=hpbInk},
      label style={font=\sffamily\fontsize{7}{8.4}\selectfont, text=hpbInk},
      xtick={1,5,10,16},
      xmajorgrids, major x grid style={draw=hpbDiscInk, line width=0.5pt,
        dash pattern=on 2pt off 2pt},
      clip=false, axis on top=false,
    },
  }
  % ------------------------------------------------------------------ (a)
  \begin{axis}[mx/common, name=wip, height=3.3cm,
      ymin=0, ymax=4.3, ytick={0,1,2,3}, xticklabels={},
      ymajorgrids, major y grid style={draw=hpbGuide, line width=0.4pt},
      ylabel={open branches $|W_\theta(t)|$},
      stack plots=y, const plot mark left,
      every axis plot/.append style={draw=hpbInk, line width=0.8pt, mark=none},
      area legend,
      reverse legend,
      legend style={draw=none, fill=white, fill opacity=0.85, text opacity=1,
        at={(0.008,0.975)}, anchor=north west, inner sep=1pt,
        legend cell align=left, row sep=-1pt,
        font=\sffamily\fontsize{7}{8.4}\selectfont, text=hpbInk},
    ]
    % stacked bottom-to-top as in the diagram: T (lane -2), X (lane -1), I (lane +1)
    \addplot[fill=hpbTask] coordinates {
      (0,0)(1,0)(2,0)(3,0)(4,0)(5,0)(6,0)(7,1)(8,1)(9,1)(10,1)(11,1)(12,1)
      (13,1)(14,1)(15,1)(16,0)(17,0)} \closedcycle;
    \addplot[fill=hpbIss] coordinates {
      (0,0)(1,0)(2,1)(3,1)(4,1)(5,0)(6,1)(7,1)(8,1)(9,1)(10,0)(11,0)(12,1)
      (13,1)(14,1)(15,1)(16,0)(17,0)} \closedcycle;
    \addplot[fill=hpbImp] coordinates {
      (0,0)(1,0)(2,0)(3,1)(4,1)(5,1)(6,1)(7,1)(8,1)(9,1)(10,0)(11,1)(12,1)
      (13,1)(14,1)(15,1)(16,0)(17,0)} \closedcycle;
    \legend{Task (T), Issue (X), Improvement (I)}
    % type letters inside the bands: the stack is readable without color
    \foreach \c/\y/\l in {12.5/0.5/T, 3/0.5/X, 12.5/1.5/X,
                          4/1.5/I, 12.5/2.5/I}{
      \edef\mxtmp{\noexpand\node[font=\noexpand\sffamily\noexpand\bfseries
        \noexpand\fontsize{7}{8.4}\noexpand\selectfont, text=hpbInk,
        inner sep=0pt] at (axis cs:\c,\y) {\l};}\mxtmp
    }
    % baseline tags above the panel
    \foreach \c/\v in {1/v1.0, 5/v1.0.1, 10/v1.1, 16/v1.2}{
      \edef\mxtmp{\noexpand\node[hpb/tag, anchor=south, yshift=3pt]
        at (axis cs:\c,4.3) {\v};}\mxtmp
    }
    % annotations
    \draw[hpb/leader] (axis cs:9.7,2.2) |- (axis cs:8.6,3.72)
      node[mx/ann, anchor=east]
      {release of v1.1: I-21 and X-102\\close together ($|C|=2$)};
    \draw[hpb/leader] (axis cs:15.7,2.2) |- (axis cs:15.1,3.72)
      node[mx/ann, anchor=east]
      {batch merge at v1.2:\\three branches close ($|C|=3$)};
  \end{axis}
  % ------------------------------------------------------------------ (b)
  \begin{axis}[mx/common, at={(wip.south west)}, anchor=north west,
      yshift=-0.55cm, height=3.6cm,
      ymin=0.4, ymax=6.6, ytick={1,...,6}, y dir=reverse,
      yticklabels={I-21, I-22, X-101, X-102, X-103, T-8},
      y tick label style={font=\sffamily\bfseries\fontsize{7}{8.4}\selectfont},
      major y tick style={draw=none},
      xticklabels={Jan, Feb, Apr, Jul},
      xlabel={ordinal time (one column per event; calendar month of each baseline)},
      ylabel={branch},
    ]
    % one bar per branch: from opening t_o(b) (start node) to closing t_c(b)
    % (release of the merge baseline);  #1 row, #2 type color, #3 letter,
    % #4 t_o column, #5 t_c column, #6 fork -> merge baseline, #7 D(b)
    \foreach \r/\tc/\tl/\o/\e/\fm/\dd in {
        1/hpbImp/I/3/10/{v1.0 $\to$ v1.1}/2,
        2/hpbImp/I/11/16/{v1.1 $\to$ v1.2}/1,
        3/hpbIss/X/2/5/{v1.0 $\to$ v1.0.1}/1,
        4/hpbIss/X/6/10/{v1.0.1 $\to$ v1.1}/1,
        5/hpbIss/X/12/16/{v1.1 $\to$ v1.2}/1,
        6/hpbTask/T/7/16/{v1.0.1 $\to$ v1.2}/2}{
      \edef\mxtmp{%
        \noexpand\draw[mx/bar, fill=\tc!55]
          (axis cs:\o,\r-0.27) rectangle (axis cs:\e,\r+0.27);
        \noexpand\node[hpb/event, fill=\tc] at (axis cs:\o,\r) {\tl};
        \noexpand\node[mx/in, anchor=west] at (axis cs:\o,\r) [xshift=9pt]
          {\unexpanded\expandafter{\fm}};
        \noexpand\node[mx/in, font=\noexpand\sffamily\noexpand\bfseries
          \noexpand\fontsize{7}{8.4}\noexpand\selectfont]
          at (axis cs:17.55,\r) {\dd};}%
      \mxtmp
    }
    \node[mx/in, font=\sffamily\bfseries\fontsize{7}{8.4}\selectfont]
      at (axis cs:17.55,0.4) [anchor=south, yshift=1pt] {$D(b)$};
  \end{axis}
  % panel letters
  \node[font=\sffamily\bfseries\fontsize{7.5}{9}\selectfont, anchor=north east]
    at ([xshift=-1.15cm]wip.north west) {(a)};
  \node[font=\sffamily\bfseries\fontsize{7.5}{9}\selectfont, anchor=north east]
    at ([xshift=-1.15cm,yshift=-0.55cm]wip.south west) {(b)};
\end{tikzpicture}}
\caption{Flow metrics read straight from the graph, for the data of
situation~A (\cref{fig:software-line}) on the same ordinal columns. (a)~Work
in progress per type: each opening raises a step, each closure lowers it, and
a release drops it by its batch size $|C(v)|$. (b)~One bar per branch, from
its opening to the release that closed it, with its fork and merge versions
and its merge distance $D(b)$. Lead times in calendar time come from the
event dates, not from the length of the bars.}
\label{fig:metrics}
\end{figure}


\paragraph{Quality and maintenance.}
The release batch is exactly the release-note list, and smaller batches go
with faster feedback \citep{reinertsen2009flow}. Escaped defects need three
cautions: \emph{seen in} is not \emph{introduced by}; counts grow with how
widely a version is deployed, so normalize by installed base where you
have it; and all of it assumes the fork baseline of every issue is the
version where it was really seen, which is why intake makes it mandatory
(\cref{subsec:method-practices}). Gate slip only works if the original
target date is kept when the date changes. Maintenance load and unported
fixes keep old references and maintenance streams honest.

\subsection{Reading metrics responsibly}
\label{subsec:metrics-goodhart}

Every metric above can be improved without improving the product. Lead
time drops if branches are opened late or split artificially. The discard
ratio drops if hopeless work is kept open. Escaped defects drop if issues
are filed against the wrong version. Gate slip disappears if scope is cut
silently. This is Goodhart's law: when a measure becomes a target, it stops
being a good measure. Three safeguards:

\begin{enumerate}[leftmargin=*,itemsep=2pt]
  \item \emph{Report in pairs} that are hard to game together: lead time
    with throughput, gate slip with descoped count, merge ratio with age
    of open work.
  \item \emph{Always show the diagram next to the numbers}, so every
    anomaly can be traced to the branches that cause it.
  \item \emph{Do not use the metrics to compare or reward people or
    teams.} Rules R1--R7 are what make the numbers trustworthy; a regime
    that rewards breaking them destroys its own data.
\end{enumerate}

\noindent The metrics describe how change flows through a product. They do
not measure customer value, effort or the quality of decisions; like
value-stream mapping \citep{rother1999lts}, they make the flow visible in
order to discuss it, not to replace the discussion.
